\documentclass{article}

\usepackage{amsmath,amssymb}
\usepackage{xurl}
\usepackage[hidelinks]{hyperref}
\setlength{\emergencystretch}{2em}

% Shared commands for notes in docs/paper-gaps/.

\DeclareUnicodeCharacter{2080}{\ifmmode{}_{0}\else\textsubscript{0}\fi}
\DeclareUnicodeCharacter{2081}{\ifmmode{}_{1}\else\textsubscript{1}\fi}
\DeclareUnicodeCharacter{2082}{\ifmmode{}_{2}\else\textsubscript{2}\fi}
\DeclareUnicodeCharacter{2099}{\ifmmode{}_{n}\else\textsubscript{n}\fi}

\newcommand{\Lean}{\textsc{Lean}}
\ExplSyntaxOn
\NewDocumentCommand{\leanid}{m}{
  \ifmmode
    \textsf{\detokenize{#1}}
  \else
    \group_begin:
    \urlstyle{sf}
    \tl_set:Nn \l_tmpa_tl {#1}
    \tl_replace_all:Nnn \l_tmpa_tl {\_} {_}
    \exp_args:NV \path \l_tmpa_tl
    \group_end:
  \fi
}
\ExplSyntaxOff
\providecommand{\tr}{\operatorname{tr}}
\newcommand{\tnleanIssueBase}{https://github.com/LionSR/TNLean/issues/}
\newcommand{\tnleanPrBase}{https://github.com/LionSR/TNLean/pull/}
\newcommand{\tnleanDocBase}{https://sirui-lu.com/TNLean/docs/}
\newcommand{\ghissue}[1]{\href{\tnleanIssueBase#1}{GitHub issue \##1}}
\newcommand{\ghpr}[1]{\href{\tnleanPrBase#1}{PR \##1}}
\newcommand{\doclink}[1]{\href{\tnleanDocBase#1.html}{\path{#1}}}

% Dirac notation and the identity operator, as in blueprint/src/macros/common.tex.
% \providecommand keeps note-local definitions authoritative.
\usepackage{dsfont}
\providecommand{\ket}[1]{|#1\rangle}
\providecommand{\bra}[1]{\langle #1|}
\providecommand{\braket}[2]{\langle #1 | #2 \rangle}
\providecommand{\ketbra}[2]{|#1\rangle\!\langle #2|}
\providecommand{\Id}{\mathds{1}}
\providecommand{\one}{\mathds{1}}
\providecommand{\C}{\mathbb{C}}
\providecommand{\MN}[1]{M_{#1}(\C)}
\providecommand{\MD}{M_D(\C)}

% External referencing for paper sources.  The build helper writes sanitized
% auxiliary files under build/paper-aux/ and excludes duplicated source labels.
\usepackage{xr}

% Import a generated paper auxiliary file with a caller-supplied label prefix.
% The candidate paths support builds from docs/paper-gaps/, the repository root,
% and build/paper-aux/ itself.
\newcommand{\importpaperaux}[2]{%
  \IfFileExists{../../build/paper-aux/#2.aux}{%
    \externaldocument[#1]{../../build/paper-aux/#2}%
  }{%
    \IfFileExists{build/paper-aux/#2.aux}{%
      \externaldocument[#1]{build/paper-aux/#2}%
    }{%
      \IfFileExists{#2.aux}{%
        \externaldocument[#1]{#2}%
      }{}%
    }%
  }%
}

% General external reference macros
\newcommand{\paperref}[2]{\ref{#1-#2}}
\newcommand{\papereqref}[2]{(\ref{#1-#2})}

% CPSV16 source (1606.00608 / MPDO-22-12-17-2.tex) external document and shortcuts
\importpaperaux{cpsv16-}{1606.00608/MPDO-22-12-17-2-tnlean}
\newcommand{\cpsvref}[1]{\paperref{cpsv16}{#1}}
\newcommand{\cpsveqref}[1]{\papereqref{cpsv16}{#1}}

% Verdict marker: \gapnote{<kind>}{<status>}. Kind is the policy
% classification (clarification, local-correction, scope-restriction,
% unfaithful, false-source, open-gap); status is open, wip (elimination
% underway), resolved, or historical. Machine-read by the site index;
% prints a verdict line.
\newcommand{\gapnote}[2]{\par\noindent\textbf{Verdict:} #1 (#2).\par}


\title{The Inverse on the Range in the Two-Variable Operator Schwarz Inequality}
\author{}
\date{2026-08-21}

\begin{document}
\maketitle

\section{Source statement}

Let $E\colon M_{d'}(\mathbb C)\to M_d(\mathbb C)$ be a two-positive linear
map.  Theorem~5.3 of Wolf's \emph{Quantum Channels \& Operations: Guided
Tour}, Chapter~5, asserts that for every $A,B\in M_{d'}(\mathbb C)$,
\begin{align}
  E(A^\dagger B)\,E(B^\dagger B)^{-1}\,E(B^\dagger A)
  &\leq E(A^\dagger A),
  \label{eq:schwarz}
\end{align}
where the inverse is taken on the range of $E(B^\dagger B)$.  The local
source is
\path{Notes/WolfNoteTexSource/ch05_schwarz_inequalities.tex}, lines
180--190.

The formal theorem proves the specialization $d'=d$ without an additional
witness hypothesis.  The inverse on the range is represented by the
Moore--Penrose inverse.  The argument below applies to arbitrary $d'$ and $d$;
the remaining restriction comes from the present definition of two-positivity,
which is stated only for endomorphisms of one matrix algebra.

\section{The range inverse}

Put
\begin{align}
  P &= E(A^\dagger A), &
  Z &= E(A^\dagger B), &
  R &= E(B^\dagger B).
\end{align}
Two-positivity gives
\begin{align}
  \begin{pmatrix}
    P & Z\\
    Z^\dagger & R
  \end{pmatrix}
  &\geq 0.
  \label{eq:block-positive}
\end{align}
The kernel clause of the Schur-complement theorem applied to
\eqref{eq:block-positive} yields
\begin{align}
  \ker R &\subseteq \ker Z.
  \label{eq:kernel-inclusion}
\end{align}
Let $R^+$ be the Moore--Penrose inverse of $R$, and let
\begin{align}
  s_R &= RR^+=R^+R
\end{align}
be the orthogonal projection onto the support of $R$.  Since $R$ is positive
semidefinite, its support is $(\ker R)^\perp$.  Equation
\eqref{eq:kernel-inclusion} is therefore equivalent to
\begin{align}
  Zs_R &= Z,
  &
  s_RZ^\dagger &= Z^\dagger.
  \label{eq:support-absorption}
\end{align}
In particular,
\begin{align}
  R\bigl(R^+Z^\dagger\bigr) &= Z^\dagger.
  \label{eq:range-witness}
\end{align}
Thus $R^+Z^\dagger v$ is the required preimage of $Z^\dagger v$ under $R$
for every vector $v$.  This is precisely the meaning of the inverse on the
range in \eqref{eq:schwarz}; no separate range--kernel duality argument is
needed.

The formal proof establishes \eqref{eq:kernel-inclusion}, then
\eqref{eq:support-absorption}, and finally
\eqref{eq:range-witness}.

\section{Derivation of the inequality}

For a fixed vector $v$, set
\begin{align}
  w &= R^+Z^\dagger v.
\end{align}
Equation~\eqref{eq:range-witness} gives $Rw=Z^\dagger v$.  The
pseudoinverse-free quadratic-form consequence of
\eqref{eq:block-positive} then gives
\begin{align}
  v^\dagger Pv
  &\geq v^\dagger Zw
   = v^\dagger ZR^+Z^\dagger v.
\end{align}
Since this holds for every $v$,
\begin{align}
  ZR^+Z^\dagger &\leq P,
\end{align}
which is \eqref{eq:schwarz}.\footnote{Formal declaration:
  \leanid{schwarz_two_variable_unconditional} in the namespace
  \leanid{SchwarzTwoVariable}.}

\section{Verdict}

Within the equal-dimensional specialization $d'=d$, there is no remaining
witness hypothesis: the phrase ``inverse on the range'' is realized by the
Moore--Penrose inverse, and the required range inclusion follows from
support-projection absorption.  The earlier proposed use of an explicit
Euclidean-space range--kernel duality is unnecessary.

Wolf's statement allows independent input and output dimensions.  Since the
formal notion of two-positivity presently applies only to
$E:M_D(\mathbb C)\to M_D(\mathbb C)$, that rectangular generality remains
unformalized.  The classification is therefore a resolved range-witness gap
together with an equal-dimensional scope restriction.

\end{document}
